\chapter{Observational strategies: power spectra and images}

\begin{bf}
  \author{Gianni Bernardi (INAF-IRA \& Rhodes University)}\\
  
Abstract\\

This chapter reviews the basics of radio interferometry (van-Citter-Zernike theorem, uv-coverage, image formation calibration) and how they are linked to the measurements of the 21-cm power spectrum and its tomography 
\end{bf}


\section{Chapter layout}

\begin{itemize}
\item review of basics of interferometry: van-Citter-Zernike theorem, $uv$-coverage, image formation;
\item from visibilities to power spectra (in particular the delay transform approach - some overlap with Trott's chapter);
\item interferometric calibration: impact on foreground subtraction, calibration errors redundant calibration;
\item array design: power spectrum vs imaging, minimmum vs maximum redundancy - some overlap with Trott's chapter?;
\item impact of calibration errors on power spectra;
\item ionospheric impact/calibration;
\item 21-cm image tomography;  
\end{itemize}

\section{Interferometry overview}

\begin{figure}[]
\begin{center}
\includegraphics[width=1.\textwidth]{Bernardi/2_element_interferometer}
%\includegraphics[trim = 0.2cm 0.6cm 0.2cm 0.2cm, scale = 1.0]{Greig/GlobalSignal}
\end{center}
%\caption{A standard schematic of the two element interferometer\footnote{Credit: http://inspirehep.net}.}
\caption{A standard schematic of the two element interferometer.}
\label{fig:fig1}
\end{figure}

The Van Cittert-Zernike theorem expresses the fundamental relationship between the sky spatial brightness (or brightness distribution) $I$ and the quantity measured by an interferometer, i.e. the visibility $V$ (e.g., \cite{TMS}):
\begin{equation}
V_{ij} ({\bf b}, \lambda) = \int_\Omega {\bar I} (\hat{\sigma}, \lambda) \, e^{-2 \pi i {\bf b} \cdot {\hat \sigma}} d {\bf \sigma},
\label{eq:1}
\end{equation}
where ${\bf b}$ is the baseline vector that separates between antenna $i$ and $j$ and $\hat{sigma}$ is the observing direction (see Figure~\ref{fig:fig1}). The baseline vector is here specified in wevelengths, i.e. ${\bf b} = \frac{{\bf b_{\rm m}}}{\lambda}$, where ${\bf b_{\rm m}}$ is the baseline vector expressed in meters and $\lambda$ is the observing wavelength. The celestial signal travels an extra path between antenna $i$ and $j$ that correspond to a geometrical time delay $\tau = {\bf b} \cdot {\hat \sigma}$, where the world ``geometrical" refers to the fact that the delay depends upon the source position in the sky and the relative separation between the two antennas.
The sky brightness distribution does not enter directly in equation~\ref{eq:1}, but filtered by the antenna primary beam response $A$ that depends upon the direction in the sky and the wavelength, i.e. ${\bar I} = A  \, I$. The response of the primary beam attenuates the sky emission away from the pointing direction, effectively reducing the field of view $\theta$ of the instrument. Generally speaking, the size of the field of view is essentially given by the antenna diameter $D$: 
\begin{equation}
\theta \approx \frac{\lambda}{D}.
\label{eq:2}
\end{equation}
The integral of equation~\ref{eq:1} is taken over the source size $\Omega$. Equation~\ref{eq:1} is often re-written in a different coordinate system, i.e. the components of the baseline vector $(u,v,w)$ and the reciprocal $(l,m,n)$, where $(l,m)$ are the coordinates in the plane on the sky tangent to the observing direction $n$ (for a detailed discussion on coordinate systems see, for example \cite{TMS}). Using this different coordinate system, equation~\ref{eq:1} becomes (e.g., \cite{TMS}):
\begin{equation}
V_{ij} (u,v,w, \lambda) = \int_\Omega {\bar I} (l, m, \lambda) \, e^{-2 \pi i (ul + vm + wn)} \frac {dl \, dm \, dn}{\sqrt{1 - l^2 - m^2}},
\label{eq:3}
\end{equation}
Although low frequency radio observations are intrinsically wide-field, for the purpose of studying the 21~cm observables, we can reduce equation~\ref{eq:3} to a two dimensional Fourier transform:
\begin{equation}
V_{ij} (u,v, \lambda) = \int_\Omega {\bar I} (l, m, \lambda) \, e^{-2 \pi i (ul + vm)} dl \, dm.
\label{eq:4}
\end{equation}
Equation~\ref{eq:4} indicates that an {\it interferometer measures the two dimensional Fourier transform of the spatial sky brightness distribution}. If our goal is to reconstruct the sky brightness distribution, equation~\ref{eq:4} can be inverted into its corresponding Fourier pair:
\begin{equation}
{\bar I} (l, m, \lambda) = \int_{- \infty}^{+ \infty} V_{ij} (u,v, \lambda) \, e^{2 \pi i (ul + vm)} du \, dv.
\label{eq:5}
\end{equation}
Equation~\ref{eq:5} is, however, a poor reconstruction of the sky brightness distribution as only one Fourier mode is sampled at the time $t$. Strictly speaking, indeed, all the quantities in equation~\ref{eq:4} and \ref{eq:5} are time variable. In most cases, the time dependence of the primary beam and the sky brightness distribution can be neglected, however, this is not the case for the visibility $V$ as the projection of the baseline vector with respect to the source direction changes significantly throughout a long (e.g. a few hours) track. In this way, many measurements of the visibility coherence function $V$ as $(u,v,)$ change with time can be made, allowing for a better reconstruction of the ${\bar I} (l, m, \lambda)$ function. This methods is commonly described as {\it filling the $uv$ plane via Earth rotation synthesis} and was invented by \cite{ryle60}. The other (complementary) way to fill the $uv$ plane is to deploy more antennas on the ground in order to increase the number of instantaneous measurements of independent Fourier modes. If $N$ antennas are connected in an interferometric array, $\frac{N (N - 1)}{2}$ instantaneous measurements are made. 

The combination of a large number of antennas and the Earth rotation synthesis, defines the sampling function $S(u,v,)$ in the $uv$ plane. In any real case, equation~\ref{eq:5} can therefore be re-written as:
\begin{equation}
{\bar I}_D (l, m, \lambda) = \int_{- \infty}^{+ \infty} S(u,v,) V (u,v, \lambda) \, e^{2 \pi i (ul + vm)} du \, dv
\label{eq:6}
\end{equation}
where ${\tilde I}_D$ indicates the sky brightness distribution sampled at a finite number of $(u,v)$ points (often termed {\it dirty image}) and I dropped the explicit dependence on the antenna pair as redundant at his point. Using the convolution theorem, equation~\ref{eq:6} can be re-written as:
\begin{equation}
{\bar I}_D (l, m, \lambda)  =  \tilde{S \, V} =  {\tilde S} \ast {\tilde V} = {\rm PSF} (l, m, \lambda) \ast {\tilde V} (l, m, \lambda),
\label{eq:6}
\end{equation}
where the tilde indicates the Fourier transform, $\ast$ the convolution operation and {\rm PSF} is the Point Spread Function, i.e. the response of the interferometric array to a point sources which, in our case, is also the Fourier transform of the $uv$ coverage.

I will give some examples of sampling functions for different instruments in Section~\ref{sec:design}, however, the sampling function always effectively reduces the integral over a finite (often not contiguous) area of the $uv$ plane. In particular, the sampled $uv$ plane is restricted to a minimum $uv$ distance that cannot be shorter than the antenna size and a maximum $uv$ distance given by equation~\ref{eq:2}. The result of not sampling all the Fourier modes is that the PSF has got ``sidelobes", i.e. nulls and secondary lobes that can often contaminate fainter true sky emission. The best reconstruction of the sky brightness distribution ${\bar I}$ requires deconvolution of the dirty image from the PSF. 




\section{21~cm observables: power spectra and images}
\label{sec:observables}

The ultimate goal of 21~cm observations is to image the spatial distribution of the 21~cm signal as a function of redshift, also known as {\it 21~cm tomography}. Given the curent theoretical predictions, such observations need to achieve mK sensitivity on a few arcminute angular scales (see Chapter 1 in this book). Most of the current arrays, however, only have the sensitivity to perform a statistical detection of the 21~cm, i.e. to measure its power spectrum. Given an intensity field $T$ function of the three dimensional spatial coordinate $\bf x$, its power spectrum $P(k)$ can be defined as:
\begin{equation}
%\langle {\tilde T}^* ({\bf k}) {\tilde T}({\bf k'}) \rangle = (2 \pi)^3 P({\bf k}) \delta^3({\bf k} - {\bf k}')
\langle {\tilde T}^* ({\bf k}) {\tilde T}({\bf k'}) \rangle = (2 \pi)^3 P(k) \delta^3({\bf k} - {\bf k}')
\label{sec_observables_eq1}
\end{equation}
where $\langle \rangle$ indicates the ensamble average, $\bf k$ is the Fourier conjugate of $\bf x$, tilde the Fourier transform, $^*$ the complex conjugate operator and $\delta$ the Dirac delta function. In 21~cm observations, power spectra can be computed directly from interferometric image cubes after deconvolution of ${\bar I}_D (l,m,\lambda)$ from the point spread function (e.g., \cite{pen09}, \cite{harker10}, \cite{beardsley16}, \cite{patil17}). An alternative way to estimate the 21~cm power spectrum that has driven the design and calibration strategies of some of the arrays. Equation~\ref{eq:4} already shows that the interferometer is a ``natural" spatial power spectrum instrument (e.g., \cite{white99}). Visibilities can be further Fourier transformed along the frequency axis (the so-called {\it delay trasform}, \cite{parsons12a}): 
\begin{equation}
{\tilde V}_{ij} (u,v, \tau) = \int_B {\bar I} (l, m, \nu) \, e^{-2 \pi i \nu \tau} d \nu
\label{sec_observables_eq2}
\end{equation}
where $B$ is the observing bandwidth and $\tau$ is the geometrical delay. The delay transform is therefore proportional to the three dimensional power spectrum (\cite{parsons12b}):
\begin{equation}
P(k) \propto {\tilde V}_{ij} (|{\bf b}|, \tau),
\label{sec_observables_eq3}
\end{equation}
where the proportionality constant transforms the visibility units into power units (\cite{parsons12b}) and the observer units $({\bf b},\tau)$ map directly in $k$ modes parallele and perpendicular to the line of sight (e.g., \cite{morales04}):
\begin{equation}
k_\perp = \frac{2 \pi |{\bf b}|}{D} = \frac{2 \pi \sqrt{u^2 + v^2}}{D}, \,\,\,\,\,\,\,    k_\parallel = \frac{2 \pi f_{21} H_0 E(z)}{c \, *1_z)^2} \, \tau,
\label{sec_observables_eq4}
\end{equation}
where $D$ is the transverse comoving distance, $f_{21} = 1421$~MHz, $H_0$ is the Hubble constant and $E(z) = \sqrt{\Omega)m (1+z)^3 + \Omega_k (1+z)^2 + \Omega_\Lambda}$. Due to the dependence of the geometrical delay upon frequency, equation~\ref{sec_observables_eq3} is only valid for short baselines, tyipically shorter than a few hundresd meters, for which the geometrical delay is fairly constant across the bandwidht and lines of constant $k_\parallel$ are practically orthogonal to the $k\perp$ axis (\cite{parsons12b}.

Equation~\ref{sec_observables_eq2} does not only provide a link between visibilities and three dimensional power spectra, but also introduces the concept of ``horizon limit" which is the maximum delay allowed 

\begin{figure}[]
\begin{center}
\includegraphics[width=1.\textwidth]{Bernardi/delay_transform}
\end{center}
\caption{A standard schematic of the two element interferometer.}
\label{fig:fig1}
\end{figure}



The power spectrum has a very close connection with the visibility measured by an interferometer (include CMB references). A direct relationship between the three dimensional 21~cm distribution and interferometric visibilities can be derived. An interferometric array already operates a two dimensional Fourier transform of the sky brightness distribution, therefore a Fourier transform 

%\section{Challenges in 21~cm observations: interferometric calibration in presence of bright foreground emission, wide fields of view, ionosphere}
\section{Interferometric calibration and 21~cm observations}
\label{sec:challenges}

Celestial signals always experience a corruption due to the non-ideal instrumental response that needs to be corrected for a posteriori, in a process that is known as interferometric calibration. Calibration relies on the definition of a data model where the corruptions are described by antenna based quantities known as Jones matrices. The data model is known as the interferometric measurement equation (\cite{hamaker96},\cite{smirnov11}) that will be summarized in the following.

If antenna $1$ and antenna $2$ measure two orthogonal, linear polarizations $x$ and $y$, the cross-polarization visibility products can be grouped in a $2 \times 2$ matrix ${\bf V}$
\begin{equation}
    {\bf V}^{12} (u,v,\lambda) \equiv 
    \left[ 
    \begin{array}{cc}
    V_{xx}^{12} (u,v,\lambda) & V_{xy}^{12} (u,v,\lambda) \\
    V_{yx}^{12} (u,v,\lambda) & V_{yy}^{12} (u,v,\lambda) \\
    \end{array}
    \right].   
\label{eq:sec:1}
\end{equation} 
The sky brightness distribution $I$ can also be written as a $2 \times 2$ matrix ${\bf B}$ using the Stokes parameters as a polarization basis:
\begin{equation}
    {\bf B}_I (l,m,\lambda) \equiv 
    \left[
    \begin{array}{cc}
    I (l,m,\lambda) + Q (l,m,\lambda) & U (l,m,\lambda) + iV (l,m,\lambda) \\
    U (l,m,\lambda) - iV (l,m,\lambda) & I (l,m,\lambda) - Q (l,m,\lambda) \\
    \end{array}
    \right].   
\label{eq:sec:2}
\end{equation} 
At this point, equation~\ref{eq:3} can be written by including the corruptions represented by the Jones matrices $J$ (\cite{hamaker96},\cite{smirnov11}):
\begin{equation}
{\bf V}^{12} (u,v,\lambda) = {\bf J}^1 \left( \int_\Omega {\bf B}_I (l, m, \lambda) \, e^{-2 \pi i (ul + vm)} dl \, dm  \right) ({\bf J}^2)^H
\label{eq:sec:3}
\end{equation} 
Equation~\ref{eq:sec:3} is known as the measurement equation and is the core of interferometric calibration. For an array with $N$ antennas, equation~\ref{eq:sec:3} can be written for each of the $\frac{N (N - 1)}{2}$ visibilities forming an overdetermined system of equations. The development of calibration algorithms is a very active research line (\cite{mitchell08}, \cite{tasse14}, \cite{smirnov15}) although beyond the scope of this chapter and we mention it here for completeness.
The solution of the system of calibration equations requires some knowledge of the sky brightness distribution ${\bf B}_I$, or a {\it sky model}. Traditionally this is achieved by observing a calibration source, i.e. a bright, unresolved point source with known spectral and polarization properties. Calibration solutions are then applied to the observed field that, in turn, is then used to improve the sky model ${\bf B}_I$ which, in turn, leads to more accurate calibration solutions $J$ in the loop that is traditionally called {\it self calibration} (\cite{cornwell81}, \cite{pearson84}). This approach can lead to a high accurate calibration (e.g., \cite{bernardi10}, \cite{smirnov11b}).

The advantage of the measurement equation is tha
t it can factorize different physical terms into different matrices. For example, the frequency response of the electronic filters and its time variations essentially affects only the two polarization response and are modeled with a diagonal Jones matrix $B$:
\begin{equation}
    {\bf B} (t,\lambda) \equiv 
    \left[
    \begin{array}{cc}
    b_x (t,\lambda) 	& 	0 	\\
    0 		& b_y (t,\lambda) 	\\
    \end{array}
    \right],   
\label{eq:sec:4}
\end{equation} 
whereas the undesired instrumental leakage between the two orthogonal polarized is represented by a $D$ Jones matrix of the form:
\begin{equation}
    {\bf D} (t,\lambda) \equiv 
    \left[
    \begin{array}{cc}
    1	 		& d_x (t,\lambda)	\\
    -d_y (t,\nu)	& 1 	\\
    \end{array}
    \right],   
\label{eq:sec:5}
\end{equation} 
and the measurement equation can be written as:
\begin{equation}
{\bf V}^{12} (u,v,\lambda) = {\bf B}^1 \, {\bf D}^1 \left( \int_\Omega {\bf B}_I (l, m, \lambda) \, e^{-2 \pi i (ul + vm)} dl \, dm  \right) ({\bf D}^2)^H  \, ({\bf B}^2)^H.
\label{eq:sec:6}
\end{equation} 
Retaining only the first order terms, equation~\ref{eq:sec:6} can be written as (\cite{sault96}):
\begin{eqnarray}
V_{xx}^{12} (u,v,\lambda) & = & b_x^1 \, b_x^{2,*} [I (l,m,\lambda) - Q (l,m,\lambda)]\\
V_{xy}^{12} (u,v,\lambda) & = & b_x^1 \, b_y^{2*} [(d_x^1 - d_y^{2*}) I (l,m,\lambda) + U (l,m,\lambda) + iV (l,m,\lambda)]	\\
V_{yx}^{12} (u,v,\lambda) & = & b_y^1 \, b_x^{2*} [(d_x^2 - d_y^{1*}) I (l,m,\lambda) + U (l,m,\lambda) - iV (l,m,\lambda)]	\\
V_{yy}^{12} (u,v,\lambda) & = & b_y^1 \, b_y^{2,*} [I (l,m,\lambda) - Q (l,m,\lambda)]
\end{eqnarray}
This simplified version of the measurement equation offers an intuitive understanding as to why calibration is so important in 21~cm observations. 


The separation of the cosmological signal from the much brighter foreground emission is the key challenge in interferometric observations of the 21~cm signal. As long as foreground emission is smooth in frequency the separation of the 21~cm signal is possible. However, it has been realized already in the early days that there is a strong coupling between foregrounds and interferometric calibration. 


\subsection{Redundant calibration}

\section{Array design}
\label{sec:design}





%Lorem ipsum dolor sit amet, consectetur adipiscing elit. Duis eu egestas erat. Maecenas tincidunt lacinia tincidunt. Mauris id lectus nec neque feugiat condimentum vitae at diam. In vel orci nunc, non commodo mauris. Vivamus ipsum enim, vulputate quis pharetra non, molestie quis felis. Vivamus porttitor placerat turpis at accumsan. Nunc tortor velit, faucibus a rhoncus nec, blandit non elit. Nam consectetur lectus eu nisi blandit dapibus rhoncus dui tempus. Mauris fermentum dolor vel ipsum vulputate sit amet ultricies tortor lacinia. Donec ut nibh erat. Morbi nec mi ante. Integer nec vestibulum diam. Donec tincidunt pellentesque quam, ut interdum mauris venenatis condimentum. Nam condimentum, augue in aliquet gravida, neque dui elementum eros, id semper eros purus sed felis. Curabitur in justo sit amet sapien ultrices hendrerit at quis nibh. Quisque iaculis pulvinar tincidunt. 
%\begin{eqnarray}
%C(12) &= &\left[\overrightarrow{\pi}\cdot\overrightarrow{\phi}(x+r)\right] \nonumber \\ 
%&\approx& 1-\mathrm{const}\frac{r^2}{L^2}\int_r^L\frac{x\rmd x}{x^2} + \cdots \nonumber  \\
%&\approx& 1-\mathrm{const}\frac{r^2}{L^2}\ln\frac{x\rmd x}{x^2} + \cdots .\label{brokenlongeqn}
%\end{eqnarray}
%
%Aenean tellus risus, porta sit amet porta vitae, tincidunt ut felis. Class aptent taciti sociosqu ad litora torquent per conubia nostra, per inceptos himenaeos. Vestibulum ante ipsum primis in faucibus orci luctus et ultrices posuere cubilia Curae; Phasellus pulvinar placerat velit auctor egestas. Vivamus euismod fringilla tincidunt. Sed ut magna felis, id sollicitudin nunc. Quisque a dui eu erat consectetur egestas a quis justo. Aenean euismod congue diam, vel posuere urna fermentum sit amet. Lorem ipsum dolor sit amet, consectetur adipiscing elit. Mauris faucibus lacus eget est mollis auctor. Donec at nibh ligula, et posuere massa. Phasellus quis leo diam \cite{diamantaras1996pcn}.
%Donec aliquam blandit risus, eu venenatis ante euismod eu. Curabitur cursus justo id arcu condimentum feugiat. Integer sapien urna, vulputate et adipiscing nec, convallis et justo. Suspendisse in ipsum at felis ornare interdum \cite{tulone2006pts},
%
%\begin{figure}[]
%\begin{center}
%\includegraphics[width=0.5\textwidth]{Bernardi/01x01-eps-converted-to}
%\end{center}
%\caption{This is figure 1 in chapter 1.}
%\end{figure}
%
%\paragraph{Cras adipiscing} sagittis nunc vel luctus. Suspendisse volutpat augue quis erat semper consequat dignissim tellus euismod. Morbi hendrerit, tellus id aliquam iaculis, nibh leo tincidunt eros, vitae varius ligula felis in mi.
%
%\begin{table}
%\caption{Greek Letters.}
%\begin{center}
%\begin{tabular}{llllllll}
%\hline
%$\alpha $  & $ \beta $  & $ \gamma $  & $ \delta $  & $ \epsilon $  & $ \varepsilon $  & $ \zeta $  & $ \eta $ \\
% $ \theta $  &  $ \vartheta $  &  $ \gamma $  &  $ \kappa $  &  $ \lambda $  &  $ \mu $  &  $ \nu $  &  $ \xi $ \\
% $ o $  &  $ \pi $  &  $ \varpi $  &  $ \rho $  &  $ \varrho $  &  $ \sigma $  &  $ \varsigma $  &  $$ \\
% $ \tau $  &  $ \upsilon $  &  $ \phi$ &  $ \varphi $  &  $ \chi $  &  $ \psi $  &  $ \omega$  &  $ $ \\
% &  &  &  &  &  &  & \\
%$ \Gamma $  & $ \Delta $  & $ \Theta $  &  $ \Lambda $  &  $ \Xi $  &  $ \Pi $  &  $ \Sigma $  & $ \Upsilon $ \\
% $ \Phi$ &  $ \Psi $  &  $ \Omega $  &  &  &  &  &\\
%\hline
%\end{tabular}
%\end{center}\end{table}
%
%\begin{figure}[]
%\begin{center}
%\includegraphics[width=0.6\textwidth]{Bernardi/01x02}
%\end{center}
%\caption{This is figure 2 in chapter 1.}
%\end{figure}
%
%
\bibliographystyle{plain}
\bibliography{Bernardi/References}


